\documentclass[11pt]{article}
\usepackage[T1]{fontenc}
\usepackage[utf8]{inputenc}
\usepackage{lmodern}
\usepackage[margin=1.02in]{geometry}
\usepackage{amsmath,amssymb,amsthm,mathtools}
\usepackage{booktabs,array,graphicx,microtype,xcolor}
\usepackage[hidelinks]{hyperref}
\hypersetup{pdftitle={The covering number C(13,7,4) is 30},pdfsubject={Covering designs and exact finite verification}}
\setlength{\parskip}{2.3pt}
\setlength{\parindent}{15pt}
\newcommand{\bin}[2]{\binom{#1}{#2}}
\newcommand{\coverfam}{\mathcal F}
\newcommand{\cat}{\mathcal C}
\newcommand{\dom}{\mathcal D}
\newcommand{\penc}{\mathcal P}
\newcommand{\XX}{X}
\newcommand{\YY}{Y}
\newcommand{\Star}{\operatorname{St}}
\newcommand{\sig}{\operatorname{sig}}
\newtheorem{theorem}{Theorem}[section]
\newtheorem{lemma}[theorem]{Lemma}
\newtheorem{proposition}[theorem]{Proposition}
\newtheorem{corollary}[theorem]{Corollary}
\theoremstyle{definition}
\newtheorem{definition}[theorem]{Definition}
\theoremstyle{remark}
\newtheorem{remark}[theorem]{Remark}
\definecolor{navy}{RGB}{24,46,67}
\title{\vspace{-1.3em}\bfseries The covering number $C(13,7,4)$ is $30$\\[0.32em]
\large\normalfont A classification-assisted exact proof}
\author{John Tripodo}
\date{29 October 2026}
\begin{document}
\maketitle
\vspace{-1.3em}
\begin{abstract}\noindent
We prove that $C(13,7,4)=30$, using the published classification of minimum
$(12,6,3)$ coverings. In a hypothetical $29$-block $(13,7,4)$ covering,
incidence counting forces at least five points of degree $15$. Fixing one of
these points leaves a classified $15$-block link and a budget of only $14$
blocks outside it. Each of four further degree-$15$ points determines a
\emph{complete outside pencil}, whose possible forms can be enumerated from
the same classification. Pairwise agreement of these pencils is necessary.
An exact census of $68\bin{12}{4}=33{,}660$ labelled root configurations leaves
$42$ cases after sound consistency filters. An independently structured,
unpruned join enumerates $280{,}152$ compatible four-pencil tuples; each uses
$15$ distinct outside blocks, contradicting the budget $14$. A concrete
$30$-block covering is verified on all $715$ quadruples. We also distinguish
this computational proof from a subsequent Lean formalization of the global
finite obstruction, whose mathematical premises are stated explicitly.
\end{abstract}
\noindent\textbf{Keywords:} covering designs; covering numbers; finite classification;
exact enumeration; computational proof; Lean formalization.
\medskip

\section{The result and its context}
For positive integers $v\geq k\geq t$, a $(v,k,t)$ \emph{covering} is a
family $\coverfam\subseteq\bin{X}{k}$, where $|X|=v$, such that every
$T\in\bin{X}{t}$ is contained in some $B\in\coverfam$. Its minimum cardinality
is the covering number $C(v,k,t)$. The problem considered here asks for the
minimum number of seven-element blocks on thirteen points needed to cover
every four-element subset.

\begin{theorem}[Main result]\label{thm:main}
\[
\boxed{C(13,7,4)=30.}
\]
\end{theorem}

An explicit construction supplies the upper bound $30$; the contribution
is the global exclusion of $29$ blocks. Prior tabulations gave the interval
$28\leq C(13,7,4)\leq30$ \cite{LJCR}, and a separately published,
formally checked elimination of one minimum-link type does not resolve
all $68$ types \cite{Dwares}. The present argument eliminates all types,
not merely one restricted instance. Related recent work on a different
covering number illustrates the increasing role of independently checked
finite certificates in design theory \cite{Krug}.

The one external mathematical ingredient is the work of Gordon,
Patashnik, Petro and Taylor \cite{GPPT}: their Theorem 4.4 proves
$C(12,6,3)=15$, and their Theorem 5.9 classifies the minimum
$(12,6,3)$ covers by a Figure~1 completion family and a single
exceptional Figure~6 cover. Their $68$ nonisomorphic minimum covers
are \emph{input}, not an original classification claimed here.
The proof below explains why that classification is enough and
isolates the precise role of the computation.

\paragraph{Proof in one paragraph.}
Every point of a hypothetical $29$-block covering occurs in at least
$15$ blocks, and there are only eight degree units above that baseline;
hence at least five points have degree exactly $15$. Root at one such
point. Its $15$ incident blocks determine a classified link, leaving
just $14$ blocks not containing the root. Each of four further minimum
points forces a full set of outside blocks. The classification supplies
exhaustive finite domains for those sets. The computation proves that
four mutually compatible choices always need at least $15$ outside
blocks. This is the contradiction.

\section{Why five minimum points are unavoidable}
Write $X=Y\cup\{\infty\}$ with $|Y|=12$ once a root is chosen.
For any point $x\in X$, define its degree and link by
\[
 d_\coverfam(x):=|\{B\in\coverfam:x\in B\}|,\qquad
 L_x(\coverfam):=\{B\setminus\{x\}:B\in\coverfam,\ x\in B\}.
\]
If $\coverfam$ covers every four-set of $X$, then $L_x(\coverfam)$ covers every
three-set of $X\setminus\{x\}$. The cited minimum-cover theorem
therefore implies $d_\coverfam(x)\geq15$ for every $x$.

\begin{lemma}[Five minimum points]\label{lem:five}
If a $(13,7,4)$ covering has exactly $29$ distinct blocks, at least
five points have degree $15$.
\end{lemma}
\begin{proof}
Double-count incidences between points and blocks:
\[
 \sum_{x\in X}d_\coverfam(x)=7|\coverfam|=7\cdot29=203
 =13\cdot15+8.
\]
Each degree exceeds $15$ by a nonnegative integer. At most eight
points can have positive excess, leaving at least five of degree $15$.
\end{proof}

Choose such a point $\infty$. The root link
$L:=L_\infty(\coverfam)$ has exactly $15$ blocks and is one of the
published minimum $(12,6,3)$ covers up to relabelling. The
\emph{outside family}
\[
 \coverfam^\circ:=\{B\in\coverfam:\infty\notin B\}
\]
contains exactly $29-15=14$ blocks, all seven-subsets of $Y$.
Choose four other degree-$15$ points $Q\subseteq Y$,
$|Q|=4$. We now show how they constrain $\coverfam^\circ$.

\section{From classified links to complete outside pencils}
Let $\cat$ be a set of representatives of the $68$ published
isomorphism classes of minimum $(12,6,3)$ coverings, each with
$15$ distinct six-subsets of a labelled twelve-element set.
Fix a labelled root link $L\in\cat$. For $q\in Y$,
the root-containing part of the $q$-link is already fixed:
\begin{equation}\label{eq:star}
 S_L(q):=\big\{\{\infty\}\cup(A\setminus\{q\}):
                 A\in L,\ q\in A\big\}.
\end{equation}
This is a family of six-subsets of $X\setminus\{q\}$.

Take a classified source link $K\in\cat$ on a twelve-point set $W$,
a source point $c\in W$, and a bijection
$\phi:W\to X\setminus\{q\}$ such that $\phi(c)=\infty$ and
\begin{equation}\label{eq:starmatch}
 \big\{\phi(A):A\in K,\ c\in A\big\}=S_L(q).
\end{equation}
Define the associated \emph{complete outside pencil}
\begin{equation}\label{eq:pencil}
 P(L,q;K,c,\phi):=
 \big\{\{q\}\cup\phi(A): A\in K,\ c\notin A\big\}
 \subseteq\bin{Y}{7}.
\end{equation}
The domain $D_L(q)$ consists of all pencils in
\eqref{eq:pencil} over all $K,c,\phi$ satisfying
\eqref{eq:starmatch}. In computation we may first prune $c$ by the
necessary equality $d_K(c)=d_L(q)$, but the definition of the
domain itself is unrestricted.

\begin{lemma}[Domain completeness]\label{lem:domain}
Suppose $\coverfam$ is a $29$-block $(13,7,4)$ covering with root link
$L$ and $d_\coverfam(q)=15$. Then
\[
 P_q:=\{B\in\coverfam^\circ:q\in B\}\in D_L(q).
\]
In particular, $P_q$ records \emph{every} outside block through $q$,
not an arbitrarily chosen subfamily.
\end{lemma}
\begin{proof}
The full link $L_q(\coverfam)$ is a minimum $(12,6,3)$ cover.
The cited classification gives a source $K\in\cat$ and a bijection
from its twelve points onto $X\setminus\{q\}$ carrying $K$ to
$L_q(\coverfam)$. If $c$ maps to $\infty$, the blocks of $K$ containing
$c$ map precisely to the known star $S_L(q)$.
The remaining blocks of $K$ are exactly the blocks of the $q$-link
that avoid $\infty$. Reattaching $q$ to their images gives
\eqref{eq:pencil}, which is exactly $P_q$.
\end{proof}

The word \emph{unrestricted} in this construction is essential.
The full map on the other eleven source points is searched, subject
to exact star equality. Restricting $\phi$ to a convenient
automorphism subgroup would not by itself justify exhaustiveness.

\paragraph{Equivariance of domain reuse.}
The implementations may reuse a domain when two known stars are
isomorphic. This does not require choosing every isomorphism between
the target stars. If $g$ is one such isomorphism, extend it by mapping
the distinguished target point to the new distinguished target point.
Composition with $g$ bijects every classified source embedding for the
first star with an embedding for the second, with inverse composition
by $g^{-1}$. Applying this bijection to all complete pencils proves
that the transported domain is the entire second domain. Equality of
coarse invariants alone is never used to authorize transport.

\section{Agreement conditions and the finite obstruction}
For distinct $q,r\in Q$ and pencils $P\in D_L(q)$,
$R\in D_L(r)$, call $(P,R)$ \emph{compatible} when
\begin{equation}\label{eq:pair}
 \{B\in P:r\in B\}=\{B\in R:q\in B\}.
\end{equation}
The two sides are the outside seven-blocks through both points.
Thus the actual pencils $P_q,P_r$ supplied by Lemma~\ref{lem:domain}
are pairwise compatible. Also, because each $P_q\subseteq\coverfam^\circ$,
\begin{equation}\label{eq:budget}
 \left|\bigcup_{q\in Q}P_q\right|\leq|\coverfam^\circ|=14.
\end{equation}
The central finite statement is the following.

\begin{proposition}[Exhaustive four-pencil obstruction]
\label{prop:finite}
For every $L\in\cat$ and every $Q\in\bin{Y}{4}$, there are
no choices $P_q\in D_L(q)$, $q\in Q$, which satisfy every pair
compatibility equation \eqref{eq:pair} and the union budget
\eqref{eq:budget}.
\end{proposition}

\begin{proof}[Exact computer-assisted verification]
The exhaustive input space contains
$68\bin{12}{4}=68\cdot495=33{,}660$ \emph{labelled} pairs $(L,Q)$.
No orbit reduction of the $Q$-sets is needed for this claim.

The computation first rejects an input if some pair of selected points
has no compatible pair of possible pencils. It then enforces arc
consistency: any pencil with no compatible partner in another
selected domain is removed, and this rule repeats until stable.
These are sound necessary filters. To see the latter formally,
consider a putative compatible four-pencil choice. Initially each
actual choice lies in its full domain. If all four remain after a
round, each is a compatible partner for each of the other three,
so none can be removed in the next round. By induction a genuine
compatible choice is never eliminated.

The recorded exact census is:
\begin{table}[ht]
\centering\small
\begin{tabular}{@{}p{0.59\linewidth}r@{}}
\toprule
Exhaustive stage & Labelled root/fourset inputs \\
\midrule
A pair has no compatible pencil choices & 31,263\\
Additional inputs closed by sound arc consistency & 2,355\\
Inputs retained for the final four-way join & 42\\
\midrule
Total: $68\bin{12}{4}$ & 33,660\\
\bottomrule
\end{tabular}
\caption{Disjoint partition of the complete Python/C++ finite search.
The last line is an exhaustive total, not an extrapolation from
representative foursets.}\label{tab:filter}
\end{table}

For each of the $42$ residual inputs, a separately written routine
constructs all six pair-compatibility relations and enumerates every
four-pencil tuple satisfying them. Importantly, this final routine
\emph{does not prune by union size}; it computes the literal union
after joining. Its result is:
\begin{table}[ht]
\centering\small
\begin{tabular}{@{}lrrr@{}}
\toprule
Root type & Residual inputs & Compatible tuples per input
 & Total tuples \\
\midrule
Figure 1 completion types & 37 & 216 & 7,992\\
Exceptional Figure 6 type & 5 & 54,432 & 272,160\\
\midrule
Total & 42 & & 280,152\\
\bottomrule
\end{tabular}
\caption{Independent unpruned join. Every one of the $280{,}152$
compatible tuples has a union of exactly $15$ distinct outside blocks.}
\label{tab:join}
\end{table}

The reconstruction of domains uses the $68$-link catalog,
with an independent route rebuilding all $28^3=21{,}952$ literal
Figure-1 completions from the published core and short blocks,
recovering $67$ completion classes under the construction action,
and separately including the exceptional Figure-6 link.
Both implementations reconstruct $68\cdot12=816$ root/point domains,
agree on the canonical digest of every domain, and obtain
Tables~\ref{tab:filter} and \ref{tab:join}.
Python ordinary and optimized executions agree; a distinct C++
implementation reproduced the decisive totals. The source, manifests,
independent joins and build receipts are itemized in
Appendix~\ref{app:reproduce}. No floating-point optimization
status, solver timeout or unverified infeasibility declaration enters
the conclusion. Every residual compatible union has $15>14$ blocks,
so \eqref{eq:budget} is impossible.
\end{proof}

\section{Proof of optimality}
\begin{proof}[Proof of Theorem~\ref{thm:main}]
Suppose a covering has at most $29$ blocks. Removing duplicate blocks,
if any, preserves coverage; we can pad with previously unused
seven-subsets to obtain a covering with exactly $29$ distinct blocks,
since $\bin{13}{7}=1716$.

By Lemma~\ref{lem:five}, choose five degree-$15$ points. Relabel one
as $\infty$ and normalize its minimum link to a representative
$L\in\cat$; this induces a relabelling of the twelve ordinary points.
Choose the other four minimum points as $Q$.
Lemma~\ref{lem:domain} supplies one complete outside pencil
$P_q\in D_L(q)$ for each $q\in Q$. All six compatibility equations
hold, and the pencils jointly occupy at most the $14$ outside blocks.
This contradicts Proposition~\ref{prop:finite}. Therefore
$C(13,7,4)\ge30$. The explicit $30$-block covering in
Appendix~\ref{app:cover}, checked against every one of the
$\bin{13}{4}=715$ four-subsets, supplies $C(13,7,4)\le30$.
\end{proof}

\section{Assurance, formalization and exact scope}
The human-readable proof relies on the \emph{published} two-part
1995 result; the computer computation checks the remaining finite
obstruction. Several independent checks address different risks:
\begin{enumerate}
\item \textbf{Source boundary.} The literal Figure-1 and Figure-6
source blocks were checked against the published article. The
$21{,}952$ completions, a $576$-element construction action, and the
$67+1$ catalog decomposition were rebuilt in a separate dependency
lock. The construction subgroup is used for the completion catalog,
not as a restriction on full star isomorphisms.
\item \textbf{Finite search.} Independently structured Python
implementations reconstruct the entire domain family and agree on
all $816$ normalized domains; the final independent four-way join
was deliberately left unpruned by the outside-block budget. The
recorded native C++ replay reproduces the decisive census.
\item \textbf{Positive and negative controls.} A valid $30$-block
cover passes all $715$ coverage tests and retains its actual minimum
point pencils when the outside budget is correctly set to $15$.
Synthetic compatible $14$-block data also survives the generic join;
malformed IDs, corrupt catalog entries and incompatible signatures
are rejected by focused tests.
\end{enumerate}

\paragraph{A distinct Lean proof.}
A later Lean package, in the private \texttt{john-lean-math}
repository, kernel-checked the global obstruction for actual covers.
Its elimination is \emph{not} the Python partition in
Table~\ref{tab:filter}. The Lean census assigns the same $33{,}660$
positions to $31{,}762$ empty-support exclusions, $1{,}893$ exact
integer-dual exclusions (including checked transports), and five
exceptional predecessor cases. These sum to $33{,}660$ but are
logically different certificate families.

The principal Lean statement
\texttt{C13.Global.exact30\_after\_published\_results} concludes that
a $30$-block cover exists and every cover has at least $30$ blocks,
under two explicit theorem parameters:
\texttt{PublishedMinimumBound} and
\texttt{PublishedClassification1995}. Another theorem,
\texttt{exact30\_after\_published\_classification}, removes the
first parameter using checked private-triple irredundance of the
catalog, but \emph{does not} establish independence from the
minimum-bound part of the historical source proof. The Lean
classification parameter itself remains unproved within Lean;
standard-axiom auditing is not a substitute for discharging a
theorem parameter.

The banked formal build receipt records a source-and-object audit of
$76{,}982$ declarations, with kernel-oriented checks rather than
\texttt{native\_decide}. The later root receipt replayed entry
modules and mutation controls against pinned cached dependencies; it
\emph{did not} recompile the entire imported dependency tree in that
later replay. A fresh full rebuild and independent mathematical
refereeing remain outstanding; no such review is claimed here.

\section{Reproducibility and limits}
The finite theorem is reproducible from the Python source, literal
classification inputs, tests, and frozen receipts included in the
\path{anc/} directory of this submission's source archive. Hexagon
preserves the source archive alongside the work. Run
\texttt{python3 anc/reproduce.py} after unpacking it. This portable
standard-library replay compares two separately structured exact
implementations, in ordinary and optimized modes, with the frozen
receipts. The native C++ cross-check described above is historical
audit evidence; its source is not included in this Python supplement.
The published classification remains an external mathematical input.

We do not claim to reprove the 1995 link classification, to have
independently established worldwide research priority, or to have
received external human peer review. Nor does this theorem classify
all optimal $30$-block covers. The precise result of the paper is the
single optimal covering number in Theorem~\ref{thm:main}.

\appendix
\section[An explicit optimal 30-block covering]{An explicit optimal $30$-block covering}\label{app:cover}
The blocks below are the published Nurmela--\"Osterg\r{a}rd
upper-bound construction tabulated by the La Jolla Covering Repository
\cite{LJCR}, not a new construction claimed here. They are seven-subsets
of $\{1,2,\ldots,13\}$. Each four-subset is contained in at least one
of the $30$ blocks. The
accompanying independent standard-library checker enumerates all
$715$ possible four-subsets, verifies distinctness and confirms
$1{,}050=30\bin{7}{4}$ total block/fourset incidences.
\begin{center}
\begingroup\footnotesize\setlength{\tabcolsep}{6pt}
\begin{tabular}{@{}ll@{}}
\toprule
Block & Block\\
\midrule
$\{1,2,5,6,7,8,9\}$ & $\{1,3,4,5,6,8,10\}$\\
$\{2,3,6,7,8,9,10\}$ & $\{2,3,4,5,8,9,11\}$\\
$\{1,3,5,7,8,10,11\}$ & $\{1,2,4,8,9,10,11\}$\\
$\{1,2,3,5,7,10,12\}$ & $\{1,3,4,5,9,10,12\}$\\
$\{1,2,4,5,6,11,12\}$ & $\{1,3,6,7,9,11,12\}$\\
$\{2,6,7,8,9,11,12\}$ & $\{4,6,7,8,9,11,12\}$\\
$\{2,3,4,6,10,11,12\}$ & $\{5,6,7,9,10,11,12\}$\\
$\{3,4,5,6,7,9,13\}$ & $\{1,2,3,5,6,10,13\}$\\
$\{1,4,6,7,9,10,13\}$ & $\{1,2,3,4,7,11,13\}$\\
$\{2,4,7,8,9,11,13\}$ & $\{2,4,5,7,10,11,13\}$\\
$\{1,3,5,9,10,11,13\}$ & $\{1,4,5,7,8,12,13\}$\\
$\{2,4,6,7,8,12,13\}$ & $\{1,2,3,8,9,12,13\}$\\
$\{3,4,7,8,10,12,13\}$ & $\{2,5,8,9,10,12,13\}$\\
$\{3,5,6,8,11,12,13\}$ & $\{2,4,6,9,11,12,13\}$\\
$\{6,7,8,9,11,12,13\}$ & $\{1,6,8,10,11,12,13\}$\\
\bottomrule
\end{tabular}
\endgroup
\end{center}
The example is different from the separate banked $30$-block cover in
\path{known30_control.json}. Both were checked independently on all $715$
four-subsets in the preparation packet.

\section{Computational protocol and frozen identifiers}\label{app:reproduce}
The accompanying supplement preserves the exact lower-bound source tree
under \path{anc/}.
Its \path{anc/reproduce.py} runs the portable checks from an unpacked
copy without access to the private repository.
The source snapshot is research-lab commit
\path{fabd230929bf17b95d6c74d738f78c553c2e0216}.
The two verifier entry points are
\path{verify_pencils.py} and \path{independent_raw_replay.py};
the independent join is in \path{peer_union_join.py}.
The frozen inputs and file hashes are recorded in
\path{minimum_link_catalog.json},
\path{domain_reconstruction_receipt.json}, and
\path{verification_manifest.json}.
The manifest identifies both program inputs and replay outputs.

From that directory, the documented portable commands are:
\begin{verbatim}
python3 verify_pencils.py > primary.json
python3 -O verify_pencils.py > primary-O.json
python3 independent_raw_replay.py > independent.json
python3 -O independent_raw_replay.py > independent-O.json
cmp primary.json primary-O.json
cmp independent.json independent-O.json
python3 -m unittest -v test_classified_pencils.py
\end{verbatim}
The historical production release gate additionally records a native
C++ cross-check. The published-dependency audit is in
\texttt{classification\_dependency\_lock.py}. The independent
joining routine must count only \emph{actually compatible} tuples;
its reported Cartesian potential is not the number of checked tuples.

\medskip\noindent
\textbf{Pinned identity data (original research snapshot):}
\begin{center}\small
\begin{tabular}{@{}p{0.41\linewidth}p{0.51\linewidth}@{}}
\toprule
Item & Frozen identifier or exact invariant\\
\midrule
Research-lab C13 checkpoint & \texttt{eec0acba559a732c98babfab497e884fc2fe9e20}\\
All root/fourset positions & $33{,}660$\\
Unpruned compatible tuples & $280{,}152$; every union has size $15$\\
Lean entry module & \path{C13Global/Final.lean}\\
Lean fixed main theorem & \path{exact30_after_published_results}\\
\bottomrule
\end{tabular}
\end{center}
The primary reconstruction domain's complete SHA-256 is the concatenation
of the following two 32-character hexadecimal strings, in order:
\begin{center}\small\ttfamily
425ec3b48e8eecdcc3eb5ac1dbe2555b\\
28420967ea7db00cb690c1e44ddbcb5e
\end{center}
The Lean entry module is in package
\path{c13-covering-lean/global-census-20261010}.
The table records existing-source evidence; it is not a claim that
the entire private lower-bound computation or Lean package was newly
re-executed for this manuscript revision. The accompanying dated audit records the fresh executions performed for
this revision, separately from historical Lean assurance.
The Lean package is not required to reproduce the computational proof
presented in this paper.

\section{Literal classification input}
For transcription checking, the Figure~1 core is the following list on
$\{0,1,\ldots,9,A,B\}$, with $A=10$ and $B=11$:
\begin{center}\ttfamily\small
014589\quad 0167AB\quad 2345AB\quad 236789\\
02468A\quad 02579B\quad 13469B\quad 13578A\\
03478B\quad 03569A\quad 12479A\quad 12568B.
\end{center}
To this core add one six-set extending each of $0123$, $4567$ and
$89AB$, choosing the other two points from the complementary eight.
The program enumerates all $28^3$ choices, including choices equivalent
under relabelling. The construction group has order $576$ and gives
$67$ representative completions. For completeness of the search it is
sufficient that every literal completion is in one of these orbits;
no claim that arbitrary cover isomorphisms preserve the core is needed.

For Figure~6 the labels are $0123456789TE$, with $T=10$ and $E=11$:
\begin{center}\ttfamily\small
012345\quad 012678\quad 0129TE\quad 0378TE\quad 0456TE\\
045789\quad 034569\quad 13689E\quad 14579E\quad 14568T\\
13457T\quad 23679T\quad 24589T\quad 24567E\quad 23458E.
\end{center}
These are transcriptions of the cited classification inputs, not a new
classification proof. Source-figure comparison and the applicability of
the published classification must be checked independently of the code.

\section{Finite-verifier logic in pseudocode}\label{app:algorithm}
The following specifies the logical search performed by the exact
finite verifier, independently of implementation-level optimizations.
Full reproducibility requires the released source and input catalog.
\begin{quote}\small\ttfamily
for each root link L in the 68-link catalog:\\
\hspace*{1em}for each Q among the 495 ordinary foursets:\\
\hspace*{2em}build each complete domain D[L,q], q in Q\\
\hspace*{2em}if some pair of domains has no compatible pair: continue\\
\hspace*{2em}repeat: discard values lacking support in another domain\\
\hspace*{2em}until no change\\
\hspace*{2em}if some domain is empty: continue\\
\hspace*{2em}enumerate all four-tuples matching all six pair relations\\
\hspace*{2em}for each compatible tuple, measure its literal block union\\
\hspace*{2em}assert union cardinality exceeds 14
\end{quote}
Arc-consistency deletion is valid by the induction in the proof of
Proposition~\ref{prop:finite}. Domain construction is exhaustive by
Lemma~\ref{lem:domain}, \emph{provided} the cited published
classification is correct and fully transcribed. The final join does
not assume a solver's answer is valid merely because it says
``infeasible.''

\begin{thebibliography}{9}
\bibitem{GPPT}
D.~M. Gordon, O.~Patashnik, J.~Petro and H.~Taylor,
\emph{Minimum $(12,6,3)$ Covers},
Ars Combinatoria \textbf{40} (1995), 161--177.
\url{https://combinatorialpress.com/ars-articles/volume-040-ars-articles/minimum-12-6-3-covers/}.

\bibitem{LJCR}
D.~M. Gordon, \emph{La Jolla Covering Repository},
$(v,k,t)=(13,7,4)$, historic upper and lower bounds.
\url{https://ljcr.dmgordon.org/cover/table.html} and
\url{https://ljcr.dmgordon.org/cover/low_tab.html}.

\bibitem{Dwares}
E.~Dwares, \emph{Formal elimination of link class 52 in the
$C(13,7,4)$ search}, public certification project (2026),
\url{https://github.com/edwares/class52-formal-certification}.
The certification concerns one class, not the full optimum.

\bibitem{Krug}
C.~Krug, \emph{The covering number $C(12,6,4)$ is 41},
preprint (2026), arXiv:2607.23766,
\url{https://arxiv.org/abs/2607.23766}.
\end{thebibliography}
\end{document}

